%%==============================================================================
%% nomenclature.tex — \input from report.lhs
%% Two tables: (1) math notation quick-reference; (2) acronyms.
%% Add a row any time you introduce a new symbol or acronym.
%%==============================================================================
\section*{Notation}
\addcontentsline{toc}{section}{Notation}

\begin{table}[h]
\centering
\begin{tabular}{ll}
\toprule
\textbf{Symbol} & \textbf{Meaning} \\
\midrule
$\cata{g}$              & catamorphism with gene $g$ \\
$\ana{g}$               & anamorphism with gene $g$ \\
$\hylo{g}{h}$           & hylomorphism \\
$\conj{f}{g}$           & split / fork $\langle f, g\rangle$ \\
$\alt{f}{g}$            & either / case analysis $[f, g]$ \\
$f \comp g$             & function composition \\
$\trans{f}$             & transpose of $f$ \\
$\kons{x}$              & constant function returning $x$ \\
$\N,\ \Z,\ \Q,\ \R,\ \B$ & naturals, integers, rationals, reals, booleans \\
$A \prodtype B$         & product type \\
$A \coprodtype B$       & coproduct (sum) type \\
$\deff$                 & "is defined as" \\
$\implied$              & proof step ("follows from") \\
\bottomrule
\end{tabular}
\caption{Notation quick-reference. Macros are defined in \texttt{ap-symbols.sty}.}
\label{tab:notation}
\end{table}

\section*{Acronyms}
\addcontentsline{toc}{section}{Acronyms}

\begin{table}[h]
\centering
\begin{tabular}{ll}
\toprule
\textbf{Acronym} & \textbf{Meaning} \\
\midrule
AoP  & Algebra of Programming \\
ADT  & Algebraic Data Type \\
FP   & Functional Programming \\
GHC  & Glasgow Haskell Compiler \\
GHCi & GHC interactive (REPL) \\
REPL & Read--Eval--Print Loop \\
\bottomrule
\end{tabular}
\caption{Acronyms used throughout this document.}
\label{tab:acronyms}
\end{table}

\newpage
